% ==========================================================================
% CUIC Quant Fund — Backtester Reference Document
% A comprehensive guide to the sports-betting backtesting framework
% ==========================================================================
\documentclass[12pt,a4paper]{report}

% --- Packages ---
\usepackage[utf8]{inputenc}
\usepackage[T1]{fontenc}
\usepackage{lmodern}
\usepackage[margin=2.5cm]{geometry}
\usepackage{amsmath,amssymb,amsthm}
\usepackage{mathtools}
\usepackage{booktabs}
\usepackage{longtable}
\usepackage{tabularx}
\usepackage{graphicx}
\usepackage{float}
\usepackage{enumitem}
\usepackage{listings}
\usepackage{xcolor}
\usepackage{hyperref}
\usepackage{cleveref}
\usepackage{algorithm}
\usepackage{algpseudocode}
\usepackage{tikz}
\usetikzlibrary{arrows.meta,positioning,shapes.geometric,calc}
\usepackage{fancyhdr}
\usepackage{titlesec}
\usepackage{tocloft}
\usepackage{parskip}
\usepackage{caption}
\usepackage{subcaption}
\usepackage{pifont}

% --- Hyperref setup ---
\hypersetup{
    colorlinks=true,
    linkcolor=blue!60!black,
    citecolor=green!50!black,
    urlcolor=blue!70!black,
    pdfauthor={James — Cardiff University Investment Club},
    pdftitle={CUIC Backtester Reference Document},
    pdfsubject={Quantitative Sports Betting Backtesting Framework},
}

% --- Code listing style ---
\definecolor{codebg}{RGB}{245,245,245}
\definecolor{codeframe}{RGB}{200,200,200}
\definecolor{codekw}{RGB}{0,0,180}
\definecolor{codecomment}{RGB}{0,128,0}
\definecolor{codestring}{RGB}{163,21,21}

\lstdefinestyle{pythonstyle}{
    language=Python,
    backgroundcolor=\color{codebg},
    frame=single,
    rulecolor=\color{codeframe},
    basicstyle=\ttfamily\small,
    keywordstyle=\color{codekw}\bfseries,
    commentstyle=\color{codecomment}\itshape,
    stringstyle=\color{codestring},
    showstringspaces=false,
    breaklines=true,
    breakatwhitespace=true,
    tabsize=4,
    numbers=left,
    numberstyle=\tiny\color{gray},
    numbersep=8pt,
    xleftmargin=1.5em,
    framexleftmargin=1em,
    captionpos=b,
}
\lstset{style=pythonstyle}

% --- Header/Footer ---
\pagestyle{fancy}
\fancyhf{}
\fancyhead[L]{\small CUIC Quant Fund}
\fancyhead[R]{\small Backtester Reference}
\fancyfoot[C]{\thepage}
\renewcommand{\headrulewidth}{0.4pt}

% --- Theorem environments ---
\theoremstyle{definition}
\newtheorem{definition}{Definition}[chapter]
\newtheorem{remark}{Remark}[chapter]
\newtheorem{example}{Example}[chapter]

% --- Custom commands ---
\newcommand{\code}[1]{\texttt{#1}}
\newcommand{\file}[1]{\texttt{#1}}
\newcommand{\func}[1]{\texttt{#1()}}
\newcommand{\pnl}{\text{PnL}}
\newcommand{\E}{\mathbb{E}}
\newcommand{\R}{\mathbb{R}}
\newcommand{\cmark}{\ding{51}}
\newcommand{\xmark}{\ding{55}}

% --- Title ---
\title{
    \vspace{-1cm}
    {\Huge\bfseries CUIC Quant Fund}\\[0.5cm]
    {\LARGE Backtester Reference Document}\\[0.3cm]
    {\large A Complete Technical Guide to the Sports-Betting\\Backtesting Framework}\\[1cm]
    {\normalsize Version 2.0 — February 2026}
}
\author{
    James\\
    Cardiff University Investment Club\\
    \texttt{james\_branch}
}
\date{26 February 2026}

% ==========================================================================
\begin{document}
% ==========================================================================

\maketitle
\thispagestyle{empty}

\vfill
\begin{center}
\textit{This document covers every module, every formula, and every design decision\\
in the CUIC backtesting framework so that no question goes unanswered.}
\end{center}
\vfill

\clearpage
\setcounter{page}{1}

% ==========================================================================
\tableofcontents
\clearpage
\listoftables
\listofalgorithms
\clearpage
% ==========================================================================


% ##########################################################################
\chapter{Introduction}
\label{ch:introduction}
% ##########################################################################

\section{Purpose of This Document}

This reference document provides a complete, self-contained explanation of the CUIC Quant Fund's sports-betting backtesting framework. It is written so that any team member—or external reviewer—can understand:

\begin{itemize}
    \item What the backtester does and why it exists,
    \item How data flows through the system from loading to final metrics,
    \item The mathematical formulas behind every metric,
    \item How walk-forward analysis prevents overfitting,
    \item How the validation and anomaly detection systems work,
    \item How to write a strategy function and integrate it.
\end{itemize}

\section{Project Context}

The Cardiff University Investment Club (CUIC) Quant Fund researches quantitative strategies for sports betting and prediction markets. The backtester is the central evaluation tool: every strategy must pass through it before being considered for live deployment.

\textbf{Key constraints:}
\begin{enumerate}
    \item \textbf{Data leakage prevention.} A strategy must never see the outcome of the game it is betting on. This is the single most critical requirement.
    \item \textbf{Standardised output.} The backtester produces an 11-column DataFrame that downstream modules (metrics, visualisation, walk-forward) depend on.
    \item \textbf{Statistical rigour.} Results must be accompanied by p-values, confidence intervals, and overfitting diagnostics.
\end{enumerate}

\section{Architecture Overview}

The framework is split into seven modules inside \file{src/cuic\_quant/}:

\begin{table}[H]
\centering
\caption{Module overview.}
\label{tab:modules}
\begin{tabularx}{\textwidth}{llX}
\toprule
\textbf{Module} & \textbf{File} & \textbf{Responsibility} \\
\midrule
Data Loader     & \file{backtest/data\_loader.py}  & Load games from Railway DB or CSV \\
Engine          & \file{backtest/engine.py}         & Core backtest loop \\
Validator       & \file{backtest/validator.py}       & 12-check result validation \\
Metrics         & \file{metrics/\_\_init\_\_.py}     & Performance metric calculations \\
Statistics      & \file{backtest/statistics.py}      & p-values, bootstrap CI, DSR, PBO \\
Walk-Forward    & \file{backtest/walk\_forward.py}   & Train/test splits, rolling/expanding windows, CPCV \\
Comparison      & \file{backtest/comparison.py}      & Multi-strategy comparison, anomaly detection \\
Kelly Criterion & \file{strategies/kelly\_criterion.py} & Optimal bet sizing \\
\bottomrule
\end{tabularx}
\end{table}

\begin{figure}[H]
\centering
\begin{tikzpicture}[
    node distance=1.2cm and 2cm,
    block/.style={rectangle, draw, rounded corners, minimum width=3cm, minimum height=1cm, align=center, fill=blue!8},
    arrow/.style={-{Stealth[length=3mm]}, thick},
]
    \node[block] (loader) {Data Loader};
    \node[block, right=of loader] (engine) {Engine\\(\func{backtest})};
    \node[block, right=of engine] (validator) {Validator\\(12 checks)};
    \node[block, below=of engine] (metrics) {Metrics\\Module};
    \node[block, below=of loader] (wf) {Walk-Forward\\Analysis};
    \node[block, below=of validator] (stats) {Statistics\\Module};
    \node[block, below=of metrics] (comp) {Comparison\\+ Anomaly};
    \node[block, above=of engine] (strategy) {Strategy\\Function};

    \draw[arrow] (loader) -- (engine);
    \draw[arrow] (strategy) -- (engine);
    \draw[arrow] (engine) -- (validator);
    \draw[arrow] (engine) -- (metrics);
    \draw[arrow] (metrics) -- (stats);
    \draw[arrow] (engine) -- (wf);
    \draw[arrow] (wf) -- (metrics);
    \draw[arrow] (metrics) -- (comp);
\end{tikzpicture}
\caption{High-level data flow through the backtesting framework.}
\label{fig:architecture}
\end{figure}


% ##########################################################################
\chapter{Data Loading}
\label{ch:data-loading}
% ##########################################################################

\section{Overview}

The function \func{load\_backtest\_data} in \file{backtest/data\_loader.py} is the entry point for all historical game data. It abstracts the data source behind a single interface.

\section{Data Sources}

\subsection{Railway PostgreSQL (Production)}

When the \code{DATABASE\_URL} environment variable is set, the loader queries the \code{sportsbook\_matches} and \code{sportsbook\_odds} tables via SQLAlchemy:

\begin{lstlisting}[caption={Database query for historical games.}]
SELECT
    m.commence_time AS timestamp,
    m.home_team || ' vs ' || m.away_team AS game,
    m.home_team, m.away_team,
    o.home_odds, o.away_odds,
    m.home_win,
    o.closing_home_odds, o.closing_away_odds
FROM sportsbook_matches m
JOIN sportsbook_odds o ON m.id = o.match_id
WHERE m.commence_time >= :start_date
  AND m.commence_time <= :end_date
ORDER BY m.commence_time ASC
\end{lstlisting}

The closing odds columns (\code{closing\_home\_odds}, \code{closing\_away\_odds}) enable Closing Line Value (CLV) computation when available.

\subsection{CSV Fallback (Development)}

If no database is configured, or the query fails, the loader reads from a local CSV file (default: \file{data/dummy\_backtest\_input.csv}).

\begin{itemize}
    \item \textbf{Strict mode} (\code{strict=True}): raises \code{RuntimeError} if the DB query fails instead of falling back.
    \item \textbf{Warning}: When fallback occurs, a \code{RuntimeWarning} is emitted so the user knows results may use synthetic data.
\end{itemize}

\section{Required Input Columns}

\begin{table}[H]
\centering
\caption{Required columns in input data.}
\label{tab:input-columns}
\begin{tabularx}{\textwidth}{llX}
\toprule
\textbf{Column} & \textbf{Type} & \textbf{Description} \\
\midrule
\code{timestamp}   & datetime  & When the game started \\
\code{game}        & str       & Human-readable match identifier (e.g., ``Lakers vs Celtics'') \\
\code{home\_team}  & str       & Home team name \\
\code{away\_team}  & str       & Away team name \\
\code{home\_odds}  & float     & Decimal odds for the home team (must be $> 1.0$) \\
\code{away\_odds}  & float     & Decimal odds for the away team (must be $> 1.0$) \\
\code{home\_win}   & int       & 1 if home team won, 0 if away team won \\
\bottomrule
\end{tabularx}
\end{table}

\textbf{Optional columns:}
\begin{itemize}
    \item \code{closing\_home\_odds}, \code{closing\_away\_odds}: Post-market closing odds for CLV computation.
\end{itemize}

\section{Sorting Guarantee}

Both the DB and CSV paths sort the output by \code{["timestamp", "game"]} to ensure deterministic ordering. This is critical because the backtest loop processes rows sequentially and chronological order prevents future data leakage.


% ##########################################################################
\chapter{The Strategy Interface}
\label{ch:strategy-interface}
% ##########################################################################

\section{What Is a Strategy Function?}

A strategy function is a Python callable that takes a game row and optional context, and returns a \emph{signal} dictionary telling the backtester what to do.

\begin{definition}[Strategy Function]
A strategy function is a callable with signature:
\[
    f: (\text{row}: \texttt{pd.Series},\ \text{context}: \texttt{dict} \mid \texttt{None}) \to \texttt{dict}
\]
\end{definition}

\section{Input: The Row}

The row is a Pandas Series containing all columns from the input data \emph{except} \code{home\_win} (and closing odds if present). This is the data leakage prevention mechanism—the strategy cannot see the outcome.

\textbf{Available fields in the row:}
\begin{itemize}
    \item \code{timestamp}, \code{game}, \code{home\_team}, \code{away\_team}
    \item \code{home\_odds}, \code{away\_odds}
\end{itemize}

\section{Input: The Context}

The context dictionary provides the strategy with historical state:

\begin{table}[H]
\centering
\caption{Context dictionary fields.}
\label{tab:context}
\begin{tabularx}{\textwidth}{llX}
\toprule
\textbf{Key} & \textbf{Type} & \textbf{Description} \\
\midrule
\code{initial\_bankroll} & \code{float}       & Starting bankroll \\
\code{bankroll}          & \code{float}       & Current bankroll \\
\code{trade\_count}      & \code{int}         & Number of trades so far \\
\code{cumulative\_pnl}   & \code{float}       & Running P\&L total \\
\code{history}           & \code{list[dict]}  & Deep copy of all past trades \\
\code{past\_games}       & \code{DataFrame}   & All previous game rows (\textbf{without} \code{home\_win}) \\
\bottomrule
\end{tabularx}
\end{table}

\begin{remark}[Deep Copy Semantics]
Both \code{history} and \code{past\_games} are deep copies. Modifying them inside the strategy function does not affect the backtester's internal state. This prevents a class of bugs where strategies accidentally corrupt shared state.
\end{remark}

\section{Output: The Signal Dictionary}

The strategy must return a dictionary with these keys:

\begin{table}[H]
\centering
\caption{Signal dictionary fields.}
\label{tab:signal}
\begin{tabularx}{\textwidth}{lllX}
\toprule
\textbf{Key} & \textbf{Type} & \textbf{Required} & \textbf{Description} \\
\midrule
\code{action}     & str   & Yes  & \code{"BUY\_HOME"}, \code{"BUY\_AWAY"}, or \code{"SKIP"} \\
\code{confidence} & float & No   & Win probability estimate $\in [0, 1]$ (used by Kelly sizing and Brier Score) \\
\code{size}       & float & No   & Dollar amount to bet (used when \code{position\_sizing=None}) \\
\code{reason}     & str   & No   & Human-readable explanation (for logging/debugging) \\
\bottomrule
\end{tabularx}
\end{table}

\section{Example Strategies}

\subsection{Always Bet Home}

The simplest possible strategy—bets \$100 on the home team for every game:

\begin{lstlisting}[caption={The always\_bet\_home reference strategy.}]
def always_bet_home(row, context=None):
    return {
        "action": "BUY_HOME",
        "confidence": 0.5,
        "size": 100.0,
        "reason": "Always bet home (test strategy)",
    }
\end{lstlisting}

This strategy exists for validation: since it is deterministic, we can verify backtest output against a known-good CSV.

\subsection{Kelly Bet Home Demo}

A demonstration of Kelly Criterion plumbing that uses implied probability + 5\% as the confidence:

\begin{lstlisting}[caption={Kelly demo strategy (DO NOT use in production).}]
def kelly_bet_home_demo(row, context=None):
    implied_prob = 1.0 / row["home_odds"]
    confidence = min(implied_prob + 0.05, 0.95)
    return {
        "action": "BUY_HOME",
        "confidence": confidence,
        "size": 100.0,
        "reason": f"DEMO kelly (implied={implied_prob:.2f})",
    }
\end{lstlisting}

\begin{remark}[Why This Is Tautological]
Since \code{confidence} $= 1/\text{odds} + 0.05$, we always have $p > 1/\text{odds}$, so Kelly will always bet. This is by construction—it demonstrates the \emph{plumbing}, not a real edge model.
\end{remark}

\section{Error Handling}

The engine handles strategy misbehaviour gracefully:

\begin{itemize}
    \item \textbf{Exception}: Strategy raises an exception $\Rightarrow$ warning emitted, row skipped, prior trades preserved.
    \item \textbf{Non-dict return}: Strategy returns \code{None}, \code{int}, \code{str}, or \code{list} $\Rightarrow$ warning, skip.
    \item \textbf{Missing \code{action} key}: $\Rightarrow$ warning with list of returned keys, treated as SKIP.
    \item \textbf{Key typos}: Common misspellings (e.g., \code{"Action"} instead of \code{"action"}) trigger a warning.
\end{itemize}


% ##########################################################################
\chapter{The Backtest Engine}
\label{ch:engine}
% ##########################################################################

\section{Overview}

The core function \func{backtest} in \file{backtest/engine.py} is the heart of the framework. It takes historical game data and a strategy function, simulates betting on each game, and returns a DataFrame of trade results.

\section{Function Signature}

\begin{lstlisting}[caption={backtest() function signature.}]
def backtest(
    data: pd.DataFrame,
    strategy_fn: Callable,
    initial_bankroll: float = 10000.0,
    cost_pct: float = 0.0,
    cost_flat: float = 0.0,
    position_sizing: str | None = None,
    kelly_fraction: float = 0.5,
) -> pd.DataFrame:
\end{lstlisting}

\section{Input Validation}
\label{sec:input-validation}

The engine fails fast on invalid parameters:

\begin{table}[H]
\centering
\caption{Input validation rules.}
\label{tab:validation-rules}
\begin{tabularx}{\textwidth}{lX}
\toprule
\textbf{Parameter} & \textbf{Rule} \\
\midrule
\code{initial\_bankroll} & Must be a finite non-negative number. NaN $\Rightarrow$ \code{ValueError}. \\
\code{cost\_pct}         & Must be in $[0, 1]$. This is a fraction, not a percentage (0.05 = 5\%). \\
\code{cost\_flat}        & Must be non-negative. \\
\code{position\_sizing}  & Case-insensitive. Only \code{"kelly"} is supported; anything else $\Rightarrow$ warning + fallback to flat. \\
\code{kelly\_fraction}   & Must be in $(0, 1]$. Only validated when \code{position\_sizing="kelly"}. \\
Required columns         & \code{timestamp}, \code{game}, \code{home\_team}, \code{away\_team}, \code{home\_odds}, \code{away\_odds}, \code{home\_win} must all be present. \\
\bottomrule
\end{tabularx}
\end{table}

\section{The Backtest Loop}
\label{sec:backtest-loop}

\begin{algorithm}[H]
\caption{Core Backtest Loop}
\label{alg:backtest}
\begin{algorithmic}[1]
\Require data, strategy\_fn, initial\_bankroll, cost\_pct, cost\_flat
\State $\text{bankroll} \gets \text{initial\_bankroll}$
\State $\text{cumulative\_pnl} \gets 0$
\State $\text{trades} \gets []$
\For{each row in data (sorted by timestamp)}
    \If{$\text{bankroll} \leq 0$}
        \State \textbf{break}  \Comment{Bankrupt — stop trading}
    \EndIf
    \State \textbf{Validate row:} skip if odds are NaN, $\leq 1.0$, or non-numeric
    \State \textbf{Validate home\_win:} skip if NaN or not in $\{0, 1\}$
    \State Build context dict with deep-copied history
    \State $\text{strategy\_row} \gets \text{row without home\_win and closing odds}$
    \State $\text{signal} \gets \text{strategy\_fn}(\text{strategy\_row}, \text{context})$
    \If{signal is not a dict, or action is SKIP}
        \State \textbf{continue}
    \EndIf
    \State Determine odds based on action (BUY\_HOME $\to$ home\_odds, BUY\_AWAY $\to$ away\_odds)
    \If{position\_sizing = ``kelly''}
        \State Apply Kelly sizing (see \Cref{sec:kelly-sizing})
    \EndIf
    \State $\text{bet\_size} \gets \min(\text{bet\_size}, \max(0, \text{bankroll} - \text{cost\_flat}))$
    \If{$\text{bet\_size} \leq 0$}
        \State \textbf{continue}
    \EndIf
    \If{won}
        \State $\pnl \gets \text{bet\_size} \times (\text{odds} - 1) \times (1 - \text{cost\_pct}) - \text{cost\_flat}$
    \Else
        \State $\pnl \gets -\text{bet\_size} - \text{cost\_flat}$
    \EndIf
    \State $\text{cumulative\_pnl} \gets \text{cumulative\_pnl} + \pnl$
    \State $\text{bankroll} \gets \text{bankroll} + \pnl$
    \State Append trade to trades list
\EndFor
\State \Return DataFrame(trades) with 11 columns
\end{algorithmic}
\end{algorithm}

\section{Data Leakage Prevention}
\label{sec:leakage-prevention}

Data leakage is the most dangerous bug in backtesting. If a strategy can see the outcome of a game before betting, it will achieve unrealistic performance. The engine prevents this at three levels:

\begin{enumerate}
    \item \textbf{Row stripping}: Before calling the strategy, the engine drops \code{home\_win}, \code{closing\_home\_odds}, and \code{closing\_away\_odds} from the row.
    \item \textbf{Past games filtering}: The \code{context["past\_games"]} DataFrame includes only rows up to (but not including) the current row, with \code{home\_win} stripped.
    \item \textbf{Deep copying}: \code{context["history"]} is a list of independent dicts (deep copy), preventing the strategy from modifying shared state.
\end{enumerate}

\section{Transaction Cost Model}
\label{sec:costs}

The engine supports two types of costs:

\begin{definition}[Transaction Costs]
For a trade with bet size $s$, odds $o$, percentage cost $c_p$, and flat cost $c_f$:
\begin{align}
    \pnl_{\text{win}}  &= s \cdot (o - 1) \cdot (1 - c_p) - c_f \label{eq:pnl-win} \\
    \pnl_{\text{loss}} &= -s - c_f \label{eq:pnl-loss}
\end{align}
\end{definition}

\begin{example}[Transaction Cost Calculation]
Bet \$100 at odds 2.50 with $c_p = 0.02$ (2\%) and $c_f = \$1.00$:
\begin{itemize}
    \item \textbf{Win}: $\pnl = 100 \times (2.50 - 1) \times (1 - 0.02) - 1.00 = 100 \times 1.50 \times 0.98 - 1.00 = 147.00 - 1.00 = \$146.00$
    \item \textbf{Loss}: $\pnl = -100 - 1.00 = -\$101.00$
\end{itemize}
\end{example}

\textbf{Bankroll capping}: The bet size is capped at $\max(0, \text{bankroll} - c_f)$ to ensure the flat fee can always be paid without the bankroll going negative.

\section{Output Columns}
\label{sec:output-columns}

\begin{table}[H]
\centering
\caption{The 11 output columns of \func{backtest}.}
\label{tab:output-columns}
\begin{tabularx}{\textwidth}{llX}
\toprule
\textbf{Column} & \textbf{Type} & \textbf{Description} \\
\midrule
\code{timestamp}       & datetime & Game start time \\
\code{game}            & str      & Match identifier \\
\code{action}          & str      & \code{"BUY\_HOME"} or \code{"BUY\_AWAY"} \\
\code{bet\_size}       & float    & Dollar amount wagered \\
\code{odds}            & float    & Decimal odds for the chosen side \\
\code{outcome}         & str      & \code{"WIN"} or \code{"LOSS"} \\
\code{pnl}             & float    & Profit or loss for this trade \\
\code{cumulative\_pnl} & float    & Running total of all P\&L \\
\code{bankroll}        & float    & Bankroll after this trade \\
\code{confidence}      & float    & Strategy's win probability estimate (NaN if not provided) \\
\code{closing\_odds}   & float    & Closing line odds for CLV computation (NaN if unavailable) \\
\bottomrule
\end{tabularx}
\end{table}

\textbf{Metadata}: The returned DataFrame stores cost parameters in \code{.attrs}:
\begin{itemize}
    \item \code{result.attrs["cost\_pct"]}
    \item \code{result.attrs["cost\_flat"]}
    \item \code{result.attrs["initial\_bankroll"]}
\end{itemize}
These are used by the validator and metrics modules to avoid parameter mismatch bugs.


% ##########################################################################
\chapter{Kelly Criterion Position Sizing}
\label{ch:kelly}
% ##########################################################################

\section{Theory}

The Kelly Criterion, introduced by John Kelly in 1956, determines the optimal fraction of bankroll to wager to maximise the expected logarithmic growth rate of wealth.

\begin{definition}[Kelly Criterion]
For a bet with win probability $p$, loss probability $q = 1 - p$, and net odds $b$ (decimal odds minus 1):
\begin{equation}
    f^* = \frac{bp - q}{b} = \frac{p(b+1) - 1}{b}
    \label{eq:kelly}
\end{equation}
where $f^*$ is the optimal fraction of bankroll to wager.
\end{definition}

\begin{remark}
If $f^* \leq 0$, the bet has no edge and should not be taken.
\end{remark}

\section{Derivation}

Starting from the expected log growth rate:
\begin{equation}
    G(f) = p \cdot \ln(1 + bf) + q \cdot \ln(1 - f)
\end{equation}

Taking the derivative and setting to zero:
\begin{align}
    \frac{dG}{df} &= \frac{pb}{1 + bf} - \frac{q}{1 - f} = 0 \\
    pb(1 - f) &= q(1 + bf) \\
    pb - pbf &= q + qbf \\
    pb - q &= f(pb + qb) = fb(p + q) = fb \\
    f^* &= \frac{pb - q}{b} = \frac{bp - (1-p)}{b}
\end{align}

\section{Fractional Kelly}

Full Kelly is mathematically optimal but produces high variance. In practice, a \emph{fractional Kelly} approach is used:

\begin{equation}
    f_{\text{actual}} = \kappa \cdot f^* \qquad \text{where } \kappa \in (0, 1]
\end{equation}

\begin{table}[H]
\centering
\caption{Common Kelly fractions and their risk profiles.}
\label{tab:kelly-fractions}
\begin{tabular}{lll}
\toprule
$\kappa$ & \textbf{Name} & \textbf{Risk Profile} \\
\midrule
0.25 & Quarter Kelly & Very conservative — minimal drawdowns \\
0.50 & Half Kelly    & \textbf{Recommended default} — good balance of growth and safety \\
0.75 & Three-quarter Kelly & Moderate — higher variance \\
1.00 & Full Kelly    & Maximum growth — high drawdown risk \\
\bottomrule
\end{tabular}
\end{table}

\begin{remark}[Why Half Kelly Is the Default]
Half Kelly sacrifices only 25\% of the growth rate compared to full Kelly, but reduces the variance by 50\%. Since the strategy's estimated win probability $p$ is never perfectly calibrated, half Kelly provides a safety margin against probability estimation errors.
\end{remark}

\section{Worked Example}

Consider a game with home odds of 1.95. Using the Kelly demo strategy:

\begin{align}
    \text{implied\_prob} &= 1 / 1.95 = 0.5128 \\
    \text{confidence} &= \min(0.5128 + 0.05, 0.95) = 0.5628 \\
    b &= 1.95 - 1 = 0.95 \\
    f^* &= \frac{0.95 \times 0.5628 - (1 - 0.5628)}{0.95} = \frac{0.5347 - 0.4372}{0.95} = \frac{0.0975}{0.95} = 0.1026 \\
    f_{\text{half}} &= 0.5 \times 0.1026 = 0.0513 \\
    \text{bet\_size} &= 0.0513 \times 10000 = \$513.16
\end{align}

\section{Implementation Details}

The engine integrates Kelly sizing as follows:

\begin{enumerate}
    \item Check if \code{position\_sizing = "kelly"} (case-insensitive).
    \item Check if the strategy returned a valid confidence value $\in (0, 1)$.
    \item If valid: compute Kelly fraction using \func{calculate\_kelly\_fraction}, multiply by current bankroll.
    \item If invalid (confidence is \code{None}, 0.0, or 1.0): fall back to the strategy's raw \code{size} field with a warning.
    \item The result is still capped at $\max(0, \text{bankroll} - c_f)$.
\end{enumerate}

\section{Maximum Fraction Cap}

To prevent catastrophic losses from a single bad probability estimate, the Kelly module enforces:
\begin{equation}
    f_{\text{final}} = \min(f_{\text{actual}},\ f_{\text{max}}) \qquad \text{default } f_{\text{max}} = 0.25
\end{equation}

This means no single bet can ever exceed 25\% of the bankroll, regardless of how high the Kelly fraction would be.


% ##########################################################################
\chapter{Performance Metrics}
\label{ch:metrics}
% ##########################################################################

All metrics are computed by \func{calculate\_all\_metrics} in \file{metrics/\_\_init\_\_.py}. The function takes a DataFrame of trades (output of \func{backtest}) and returns a dictionary of computed values.

\section{Core Metrics}

\subsection{Total Trades}
\begin{equation}
    N = |\text{trades}|
\end{equation}
Simply the number of executed trades (rows in the output DataFrame).

\subsection{Win Rate}
\begin{equation}
    W = \frac{\text{number of WIN outcomes}}{N}
\end{equation}

\subsection{Total P\&L}
\begin{equation}
    \text{Total PnL} = \sum_{i=1}^{N} \pnl_i
\end{equation}

\subsection{Profit Factor}
\begin{equation}
    \text{Profit Factor} = \frac{\sum_{i : \pnl_i > 0} \pnl_i}{\left|\sum_{i : \pnl_i < 0} \pnl_i\right|}
    \label{eq:profit-factor}
\end{equation}

A profit factor $> 1$ means the strategy is profitable. Returns $\infty$ if there are no losses, and $0$ if there are no wins.

\section{Risk-Adjusted Metrics}

\subsection{Percentage Returns}

Before computing Sharpe and Sortino, we convert P\&L to percentage returns:
\begin{equation}
    r_i = \frac{\pnl_i}{B_{i-1}}
    \label{eq:pct-return}
\end{equation}
where $B_{i-1}$ is the bankroll \emph{before} trade $i$ (computed as $B_i - \pnl_i$ from the output).

\subsection{Annualisation Factor}
\label{sec:annualisation}

The annualisation factor is computed from actual bet frequency, not hardcoded:

\begin{equation}
    P = \frac{N}{\Delta T / 365.25}
\end{equation}

where $\Delta T$ is the time span in days between the first and last trade. If timestamps are unavailable, falls back to 365 with a warning.

\subsection{Sharpe Ratio}
\label{sec:sharpe}

\begin{definition}[Annualised Sharpe Ratio]
\begin{equation}
    \text{SR} = \frac{\bar{r} - r_f / P}{\sigma_r} \times \sqrt{P}
    \label{eq:sharpe}
\end{equation}
where $\bar{r}$ is the mean excess return, $\sigma_r$ is the standard deviation of returns (using Bessel's correction, \code{ddof=1}), $r_f$ is the annual risk-free rate, and $P$ is the annualisation factor.
\end{definition}

Returns 0.0 if there is only one data point or zero variance.

\subsection{Sortino Ratio}
\label{sec:sortino}

\begin{definition}[Annualised Sortino Ratio]
\begin{equation}
    \text{Sortino} = \frac{\bar{r} - r_f / P}{\sigma_d} \times \sqrt{P}
    \label{eq:sortino}
\end{equation}
where $\sigma_d$ is the \emph{downside deviation}:
\begin{equation}
    \sigma_d = \sqrt{\frac{1}{N} \sum_{i=1}^{N} \left[\min(0, r_i - r_f/P)\right]^2}
    \label{eq:downside-dev}
\end{equation}
\end{definition}

\begin{remark}[Sharpe vs Sortino]
The Sharpe ratio penalises all volatility equally—including upside volatility, which is desirable. The Sortino ratio only penalises \emph{downside} volatility, making it more appropriate for asymmetric return distributions (common in sports betting where $\pnl_{\text{win}} \neq |\pnl_{\text{loss}}|$).
\end{remark}

\begin{remark}[Downside Deviation Formula]
Note that the downside deviation uses \textbf{all $N$ observations} in the denominator, not just the negative ones. This is the standard definition from Sortino \& Price (1994). Using only negative observations would overstate the downside risk.
\end{remark}

\subsection{Maximum Drawdown}
\label{sec:max-drawdown}

\begin{definition}[Maximum Drawdown]
\begin{equation}
    \text{MDD} = \max_{t} \left(\frac{\text{peak}_t - E_t}{\text{peak}_t}\right)
    \label{eq:max-drawdown}
\end{equation}
where $E_t = B_0 + \sum_{i=1}^{t} \pnl_i$ is the equity at time $t$, and $\text{peak}_t = \max_{s \leq t} E_s$ is the running maximum equity.
\end{definition}

The equity curve is anchored to \code{initial\_bankroll} ($B_0$), so early losses are measured relative to the starting equity, not zero.

\subsection{Calmar Ratio}

\begin{equation}
    \text{Calmar} = \frac{R_{\text{annual}}}{\text{MDD}}
    \label{eq:calmar}
\end{equation}

where $R_{\text{annual}}$ is the annualised return (total return divided by time span in years). Returns 0.0 if MDD is zero.

\section{Betting-Specific Metrics}

\subsection{Return on Investment (ROI)}
\begin{equation}
    \text{ROI} = \frac{\text{Total PnL}}{\text{Total Wagered}} = \frac{\sum \pnl_i}{\sum s_i}
\end{equation}

This is \emph{yield on turnover}—how much profit you make per dollar wagered.

\subsection{Return on Capital}
\begin{equation}
    \text{ROC} = \frac{\text{Total PnL}}{B_0}
\end{equation}

This is the ``true'' ROI—how much profit relative to starting bankroll.

\subsection{Yield Per Bet}
\begin{equation}
    \text{Yield/Bet} = \frac{1}{N_{\text{valid}}} \sum_{i : s_i > 0} \frac{\pnl_i}{s_i}
    \label{eq:yield-per-bet}
\end{equation}

\begin{remark}[Mean of Ratios vs Ratio of Means]
This uses the \textbf{mean of per-trade ratios}, not the ratio of the means ($\text{Total PnL} / \text{Total Wagered}$). The mean-of-ratios approach gives equal weight to each trade regardless of size, which is more informative when bet sizes vary significantly (e.g., under Kelly sizing).
\end{remark}

\subsection{Average Odds}
\begin{equation}
    \bar{o} = \frac{1}{N} \sum_{i=1}^{N} o_i
\end{equation}

\subsection{Bet Frequency}
\begin{equation}
    \text{Freq} = \frac{N}{\Delta T_{\text{days}}}
\end{equation}
Bets per day, computed from the timestamp span.

\subsection{Kelly Growth Rate}
\label{sec:kelly-growth}

\begin{equation}
    G = \frac{1}{N} \sum_{i=1}^{N} \ln\!\left(1 + \frac{\pnl_i}{B_{i-1}}\right)
    \label{eq:kelly-growth}
\end{equation}

The \emph{realised} geometric growth rate per bet. The Kelly Criterion maximises the \emph{expected} value of this quantity. Positive means the bankroll is growing; negative means it is shrinking.

Returns are clipped at $-0.9999$ to prevent $\ln(0)$ when the bankroll is nearly wiped out.

\subsection{Closing Line Value (CLV)}
\label{sec:clv}

\begin{definition}[Closing Line Value]
\begin{equation}
    \text{CLV} = \frac{1}{N_{\text{valid}}} \sum_{i=1}^{N_{\text{valid}}} \left(\frac{o_i^{\text{bet}}}{o_i^{\text{close}}} - 1\right)
    \label{eq:clv}
\end{equation}
where $o_i^{\text{bet}}$ is the odds at which the bet was placed and $o_i^{\text{close}}$ is the closing line odds.
\end{definition}

CLV is the \textbf{gold standard metric in sports betting}. A positive CLV means you consistently got better prices than the final market. CLV detects genuine skill in approximately 50 bets, whereas P\&L requires 2000+ bets for statistical significance.

\section{Probability Calibration Metrics}

\subsection{Brier Score}
\label{sec:brier}

\begin{equation}
    \text{BS} = \frac{1}{N} \sum_{i=1}^{N} (p_i - y_i)^2
    \label{eq:brier}
\end{equation}

where $p_i$ is the strategy's confidence and $y_i \in \{0, 1\}$ is the actual outcome (1 = WIN).

\begin{itemize}
    \item $\text{BS} = 0$: perfect calibration
    \item $\text{BS} = 0.25$: equivalent to always predicting $p = 0.5$
    \item $\text{BS} = 1.0$: worst possible
\end{itemize}

Confidence values outside $[0, 1]$ are clamped with a warning.

\subsection{Log Loss (Cross-Entropy)}
\label{sec:logloss}

\begin{equation}
    \text{LL} = -\frac{1}{N} \sum_{i=1}^{N} \left[y_i \ln(p_i) + (1 - y_i) \ln(1 - p_i)\right]
    \label{eq:logloss}
\end{equation}

Log loss penalises confident wrong predictions more heavily than Brier score. Probabilities are clipped to $[10^{-15}, 1 - 10^{-15}]$ to avoid $\ln(0)$.


% ##########################################################################
\chapter{Result Validation}
\label{ch:validation}
% ##########################################################################

The validator (\file{backtest/validator.py}) runs 12 checks across three categories to ensure backtest results are correct and free of data leakage.

\section{Validation API}

\begin{lstlisting}[caption={Validator function signature.}]
def validate_backtest_results(
    results: pd.DataFrame,
    input_data: pd.DataFrame,
    initial_bankroll: float | None = None,
    cost_pct: float | None = None,
    cost_flat: float | None = None,
) -> dict:
\end{lstlisting}

When \code{initial\_bankroll}, \code{cost\_pct}, or \code{cost\_flat} are \code{None}, the validator auto-reads them from \code{results.attrs}. This eliminates parameter mismatch bugs where the validator uses different parameters than the backtest.

\section{Category 1: Schema Validation (5 Checks)}

\begin{table}[H]
\centering
\caption{Schema validation checks.}
\label{tab:schema-checks}
\begin{tabularx}{\textwidth}{clX}
\toprule
\textbf{\#} & \textbf{Check} & \textbf{What It Verifies} \\
\midrule
1 & Column names    & All 11 required columns are present (extra columns like \code{closing\_odds} are OK) \\
2 & Valid actions   & Every action is \code{BUY\_HOME} or \code{BUY\_AWAY} (no \code{SKIP}—those don't appear in output) \\
3 & Valid outcomes  & Every outcome is \code{WIN} or \code{LOSS} \\
4 & Positive bets   & No bet\_size is zero or negative \\
5 & Valid odds      & All decimal odds are $> 1.0$ \\
\bottomrule
\end{tabularx}
\end{table}

\textbf{Critical}: If check 1 fails (missing columns), the validator returns immediately since subsequent checks depend on those columns.

\section{Category 2: Math Correctness (4 Checks)}

\begin{table}[H]
\centering
\caption{Math correctness checks.}
\label{tab:math-checks}
\begin{tabularx}{\textwidth}{clX}
\toprule
\textbf{\#} & \textbf{Check} & \textbf{What It Verifies} \\
\midrule
6 & PnL formula & WIN: $\pnl = s(o-1)(1-c_p) - c_f$; LOSS: $\pnl = -s - c_f$. Tolerance: $\pm\$0.01$. \\
7 & Cumulative PnL & $\text{cumulative\_pnl}_i = \sum_{j=1}^{i} \pnl_j$. Verified as running sum. \\
8 & Bankroll tracking & $\text{bankroll}_i = B_0 + \text{cumulative\_pnl}_i$. \\
9 & No overbetting & $s_i \leq B_{i-1} - c_f$ (bet never exceeds bankroll minus flat cost). \\
\bottomrule
\end{tabularx}
\end{table}

\section{Category 3: Data Leakage Detection (3 Checks)}

\begin{table}[H]
\centering
\caption{Data leakage detection checks.}
\label{tab:leakage-checks}
\begin{tabularx}{\textwidth}{clX}
\toprule
\textbf{\#} & \textbf{Check} & \textbf{What It Verifies} \\
\midrule
10 & Game existence   & Every game in results exists in the original input data \\
11 & Outcome consistency & WIN/LOSS in results matches \code{home\_win} in input (e.g., BUY\_HOME + home\_win=1 $\to$ WIN) \\
12 & Chronological order & Trades are in timestamp order (out-of-order could indicate future data access) \\
\bottomrule
\end{tabularx}
\end{table}

\section{Return Value}

\begin{lstlisting}[caption={Validation return structure.}]
{
    "passed": True,          # All checks passed
    "checks_run": 12,        # Total checks executed
    "checks_passed": 12,     # How many passed
    "failures": [],          # List of failure messages
}
\end{lstlisting}


% ##########################################################################
\chapter{Statistical Significance Testing}
\label{ch:statistics}
% ##########################################################################

The statistics module (\file{backtest/statistics.py}) provides tools to determine whether backtest results are statistically meaningful or could be due to chance.

\section{Binomial Test P-Value}
\label{sec:p-value}

\begin{definition}[Binomial Test]
Given an observed win rate $\hat{p}$ over $n$ trades, the p-value tests:
\begin{equation}
    H_0: p \leq p_0 \qquad \text{vs} \qquad H_1: p > p_0
    \label{eq:binomial-test}
\end{equation}
where $p_0 = 0.5$ (no edge) by default.
\end{definition}

The test uses a \textbf{one-sided ``greater''} alternative by default, because in betting we only care if the strategy is \emph{better} than the null—a 30\% win rate is not ``significant'' in a useful sense.

The exact binomial test (via \code{scipy.stats.binomtest}) is used rather than a normal approximation.

\section{Bootstrap Confidence Intervals}
\label{sec:bootstrap}

\begin{algorithm}[H]
\caption{Bootstrap Confidence Interval}
\begin{algorithmic}[1]
\Require PnL series, $B = 10000$ resamples, confidence level $\alpha = 0.95$
\State $\hat{\theta} \gets \text{metric\_fn}(\text{data})$ \Comment{Point estimate}
\For{$b = 1$ to $B$}
    \State $\text{sample}_b \gets$ resample data with replacement
    \State $\theta_b^* \gets \text{metric\_fn}(\text{sample}_b)$
\EndFor
\State $\text{lower} \gets P_{(1-\alpha)/2}(\theta_1^*, \ldots, \theta_B^*)$
\State $\text{upper} \gets P_{(1+\alpha)/2}(\theta_1^*, \ldots, \theta_B^*)$
\State \Return $(\text{lower}, \hat{\theta}, \text{upper})$
\end{algorithmic}
\end{algorithm}

The bootstrap is non-parametric—it makes no assumptions about the distribution of returns, which is important for sports betting where returns are highly non-normal.

\section{Minimum Sample Size}
\label{sec:min-sample}

Using the normal approximation to the binomial:

\begin{equation}
    n = \frac{\left(z_{\alpha/2}\sqrt{p_0(1-p_0)} + z_{\beta}\sqrt{p_1(1-p_1)}\right)^2}{(p_1 - p_0)^2}
    \label{eq:min-sample}
\end{equation}

where $p_0 = 0.5$, $p_1 = p_0 + \text{edge}$, $z_{\alpha/2}$ is the critical value for significance $\alpha$, and $z_{\beta}$ is the critical value for power $1 - \beta$.

\begin{example}
For a 5\% edge ($p_1 = 0.55$), power = 0.80, $\alpha = 0.01$:
\[
    n = \frac{(2.576 \times 0.500 + 0.842 \times 0.497)^2}{0.05^2} \approx 1070 \text{ bets}
\]
\end{example}

\section{Significance Report}

The function \func{significance\_report} combines all of the above:
\begin{itemize}
    \item P-value from binomial test
    \item Bootstrap 95\% confidence intervals for mean PnL and ROI
    \item Sample size adequacy assessment
    \item Warnings when the sample is too small
\end{itemize}


% ##########################################################################
\chapter{Overfitting Protection}
\label{ch:overfitting}
% ##########################################################################

When testing multiple strategies on the same data, some will appear profitable purely by chance. This chapter covers the tools that detect and correct for this.

\section{The Multiple Testing Problem}

If you test $k$ strategies at significance level $\alpha$, the probability of at least one false positive is:
\begin{equation}
    P(\text{at least one false positive}) = 1 - (1 - \alpha)^k
\end{equation}

For $k = 55$ (11 team members $\times$ 5 strategies each) and $\alpha = 0.05$:
\[
    P = 1 - 0.95^{55} = 0.942
\]

There is a 94.2\% chance of finding a ``significant'' result that is actually noise.

\section{Bonferroni Correction}
\label{sec:bonferroni}

\begin{definition}[Bonferroni Correction]
Reject $H_0$ for strategy $i$ if:
\begin{equation}
    p_i < \frac{\alpha}{k}
    \label{eq:bonferroni}
\end{equation}
where $k$ is the total number of strategies tested.
\end{definition}

This is the simplest and most conservative correction. For $k = 55$ and $\alpha = 0.05$, the threshold becomes $0.05 / 55 = 0.00091$.

\section{Holm-Bonferroni Step-Down Procedure}
\label{sec:holm}

Less conservative than Bonferroni. It sorts p-values in ascending order and rejects while:
\begin{equation}
    p_{(i)} < \frac{\alpha}{k - i + 1}
\end{equation}

Once a test fails to reject, all subsequent tests also fail.

\section{Deflated Sharpe Ratio}
\label{sec:dsr}

\begin{definition}[Deflated Sharpe Ratio (Bailey \& Lopez de Prado, 2014)]
The DSR tests whether the observed Sharpe ratio exceeds what we'd expect from the best of $k$ random strategies:

\begin{equation}
    \text{DSR} = 1 - \Phi\!\left(\frac{\text{SR}_{\text{obs}} - \text{SR}_{\text{max}}^*}{\text{SE}(\text{SR})}\right)
    \label{eq:dsr}
\end{equation}
\end{definition}

The expected maximum Sharpe from $k$ random strategies uses the Euler-Mascheroni approximation:
\begin{equation}
    \text{SR}_{\text{max}}^* = \sigma_{\text{SR}} \left[(1 - \gamma)\Phi^{-1}\!\left(1 - \frac{1}{k}\right) + \gamma\,\Phi^{-1}\!\left(1 - \frac{1}{ke}\right)\right]
    \label{eq:expected-max-sharpe}
\end{equation}
where $\gamma \approx 0.5772$ is the Euler-Mascheroni constant.

The standard error of the Sharpe ratio uses the full formula that accounts for non-normality:
\begin{equation}
    \text{SE}(\text{SR}) = \sqrt{\frac{1 - \hat{\gamma}_3 \cdot \text{SR} + \frac{\hat{\gamma}_4 - 3}{4} \cdot \text{SR}^2}{n - 1}}
    \label{eq:se-sharpe}
\end{equation}
where $\hat{\gamma}_3$ is the skewness and $\hat{\gamma}_4$ is the kurtosis of returns.

\begin{remark}[Why Skewness and Kurtosis Matter]
Binary betting returns are typically highly skewed ($\hat{\gamma}_3 \in [-3, -1]$) with excess kurtosis ($\hat{\gamma}_4 > 5$). Ignoring these terms (i.e., assuming Gaussian returns) understates the standard error and makes the DSR too lenient, leading to false positives.
\end{remark}

\section{Probability of Backtest Overfitting (PBO)}
\label{sec:pbo}

\begin{definition}[Simplified PBO]
Given in-sample Sharpe ratios $\{\text{SR}_i^{\text{IS}}\}$ and out-of-sample Sharpe ratios $\{\text{SR}_i^{\text{OOS}}\}$ for $k$ strategies:

\begin{enumerate}
    \item Find the strategy with the highest IS Sharpe: $i^* = \arg\max_i \text{SR}_i^{\text{IS}}$
    \item Count how many other strategies have higher OOS Sharpe: $m = |\{j \neq i^* : \text{SR}_j^{\text{OOS}} > \text{SR}_{i^*}^{\text{OOS}}\}|$
    \item $\text{PBO} = \frac{m}{k - 1}$
\end{enumerate}
\end{definition}

\begin{remark}
This is a simplified single-split rank test, NOT the full Combinatorial Symmetric Cross-Validation (CSCV) algorithm from Bailey et al.\ (2017). For rigorous PBO, use \func{combinatorial\_purged\_cv} with multiple split combinations.
\end{remark}

\section{Overfitting Report}

The function \func{overfitting\_report} combines:
\begin{itemize}
    \item Bonferroni correction results
    \item Holm-Bonferroni correction results
    \item Deflated Sharpe Ratio p-values for each strategy
    \item Expected false positive count: $k \times \alpha$
    \item Warnings when the number of significant strategies is close to the expected false positive count
\end{itemize}


% ##########################################################################
\chapter{Walk-Forward Analysis}
\label{ch:walk-forward}
% ##########################################################################

Walk-forward analysis is the most important technique for producing realistic out-of-sample performance estimates. Fitting a model on data and evaluating on the same data produces meaningless results.

\section{Simple Train/Test Split}
\label{sec:train-test}

\begin{equation}
    \text{data} = \underbrace{[\text{row}_1, \ldots, \text{row}_k]}_{\text{train}} \cup \underbrace{[\text{row}_{k+g+1}, \ldots, \text{row}_n]}_{\text{test}}
\end{equation}

where $k = \lfloor n \times r \rfloor$ (train ratio), and $g$ is the purge gap.

The gap prevents information leakage when strategies use features with look-back windows. If a feature uses the last 5 games, setting $g = 5$ ensures no test-period data contaminates training features.

\begin{lstlisting}[caption={Train/test split usage.}]
train, test = train_test_split(data, train_ratio=0.7, gap=5)
\end{lstlisting}

\textbf{Chronological check}: If the data has a \code{timestamp} column and is not sorted in ascending order, a warning is emitted.

\section{Rolling Window Walk-Forward}
\label{sec:rolling-wf}

\begin{figure}[H]
\centering
\begin{tikzpicture}[
    scale=0.8,
    train/.style={fill=blue!20, draw=blue!50, minimum height=0.6cm},
    test/.style={fill=red!20, draw=red!50, minimum height=0.6cm},
]
    % Timeline
    \draw[->] (0,0) -- (14,0) node[right] {time};

    % Fold 0 (skipped)
    \node[test, minimum width=3cm] at (1.5, 2.5) {};
    \node at (-1, 2.5) {\footnotesize Fold 0 (skipped)};
    \node at (1.5, 2.5) {\footnotesize test};
    \node at (6, 2.5) {\footnotesize \textit{zero training data}};

    % Fold 1
    \node[train, minimum width=3cm] at (1.5, 1.5) {};
    \node[test, minimum width=3cm] at (5, 1.5) {};
    \node at (-1, 1.5) {\footnotesize Fold 1};
    \node at (1.5, 1.5) {\footnotesize train};
    \node at (5, 1.5) {\footnotesize test};

    % Fold 2
    \node[train, minimum width=5cm] at (3.5, 0.5) {};
    \node[test, minimum width=3cm] at (8.5, 0.5) {};
    \node at (-1, 0.5) {\footnotesize Fold 2};
    \node at (3.5, 0.5) {\footnotesize train};
    \node at (8.5, 0.5) {\footnotesize test};
\end{tikzpicture}
\caption{Rolling window walk-forward with 3 splits. Fold 0 is skipped because it has zero training data.}
\label{fig:rolling-wf}
\end{figure}

The data is divided into $n$ equal-sized test windows. For each fold $i$:
\begin{align}
    \text{test\_start}_i &= i \times \lfloor N / n \rfloor \\
    \text{train\_start}_i &= \max(0, \text{test\_start}_i - g - T) \\
    \text{train\_end}_i &= \max(\text{train\_start}_i, \text{test\_start}_i - g)
\end{align}

where $T = \lfloor \text{fold\_size} \times r / (1-r) \rfloor$ is the training window size.

\textbf{Fold 0 skip}: The first fold typically has zero training data (no history to train on). It is automatically skipped with a warning.

\section{Expanding Window Walk-Forward}
\label{sec:expanding-wf}

Unlike rolling window (where the training set has a fixed size), expanding window grows the training set over time:

\begin{figure}[H]
\centering
\begin{tikzpicture}[
    scale=0.8,
    train/.style={fill=blue!20, draw=blue!50, minimum height=0.6cm},
    test/.style={fill=red!20, draw=red!50, minimum height=0.6cm},
]
    \draw[->] (0,0) -- (14,0) node[right] {time};

    % Fold 0
    \node[train, minimum width=5cm] at (2.5, 2) {};
    \node[test, minimum width=2cm] at (6, 2) {};
    \node at (-1, 2) {\footnotesize Fold 0};
    \node at (2.5, 2) {\footnotesize train (50)};
    \node at (6, 2) {\footnotesize test};

    % Fold 1
    \node[train, minimum width=7cm] at (3.5, 1) {};
    \node[test, minimum width=2cm] at (8, 1) {};
    \node at (-1, 1) {\footnotesize Fold 1};
    \node at (3.5, 1) {\footnotesize train (60)};
    \node at (8, 1) {\footnotesize test};

    % Fold 2
    \node[train, minimum width=9cm] at (4.5, 0) {};
    \node[test, minimum width=2cm] at (10, 0) {};
    \node at (-1, 0) {\footnotesize Fold 2};
    \node at (4.5, 0) {\footnotesize train (70)};
    \node at (10, 0) {\footnotesize test};
\end{tikzpicture}
\caption{Expanding window walk-forward. Training set grows with each fold.}
\label{fig:expanding-wf}
\end{figure}

Parameters:
\begin{itemize}
    \item \code{min\_train\_size}: minimum rows before the first test (default 50)
    \item \code{step\_size}: rows added per fold (default 10)
\end{itemize}

\section{Anchored Walk-Forward}
\label{sec:anchored-wf}

Training always begins at a fixed \code{anchor\_date} and extends up to the start of each test period. Useful when you want a specific historical event always included in training.

\begin{lstlisting}[caption={Anchored walk-forward usage.}]
results = anchored_walk_forward(
    data, strategy_fn,
    anchor_date="2025-01-01",
    test_periods=[
        ("2025-04-01", "2025-05-01"),
        ("2025-05-01", "2025-06-01"),
    ],
)
\end{lstlisting}

\section{Combinatorial Purged Cross-Validation (CPCV)}
\label{sec:cpcv}

\begin{definition}[CPCV (Lopez de Prado)]
Data is divided into $n$ contiguous groups. All $\binom{n}{k}$ combinations of $k$ groups are used as test sets. The remaining $n - k$ groups form the training set. Rows within a \code{purge\_gap} of a train/test boundary are removed from training to prevent leakage.
\end{definition}

For $n = 5$ splits and $k = 2$ test splits: $\binom{5}{2} = 10$ combinations.

CPCV provides the most comprehensive overfitting assessment because it tests every possible train/test partition, not just sequential ones.

\section{Metric Aggregation Across Folds}
\label{sec:metric-aggregation}

Different metrics require different aggregation methods:

\begin{table}[H]
\centering
\caption{How metrics are aggregated across walk-forward folds.}
\label{tab:aggregation}
\begin{tabularx}{\textwidth}{lX}
\toprule
\textbf{Metric} & \textbf{Aggregation Method} \\
\midrule
Total PnL       & Sum across folds \\
Total trades     & Sum across folds \\
Win rate         & Trade-weighted average: $\sum_i W_i N_i / \sum_i N_i$ \\
Profit factor    & Trade-weighted average \\
Sharpe ratio     & \textbf{Concatenate} all OOS trades, compute single Sharpe on combined series \\
Sortino ratio    & \textbf{Concatenate} all OOS trades, compute single Sortino on combined series \\
Max drawdown     & Maximum (worst) across all folds \\
\bottomrule
\end{tabularx}
\end{table}

\begin{remark}[Why Concatenation for Sharpe/Sortino]
Averaging Sharpe ratios across folds is mathematically invalid because Sharpe is not a linear function—$\text{mean}(\text{SR}_1, \text{SR}_2) \neq \text{SR}(\text{combined})$. The correct approach is to concatenate all out-of-sample trade results into a single series and compute one Sharpe/Sortino ratio.
\end{remark}

\section{Walk-Forward Report}

The function \func{walk\_forward\_report} produces a human-readable multi-line report:
\begin{itemize}
    \item Aggregated OOS metrics (total trades, PnL, win rate, Sharpe, Sortino, MDD, profit factor)
    \item Per-fold breakdown (train/test sizes, OOS trades, PnL, win rate)
    \item In-sample vs out-of-sample comparison (PnL and Sharpe for each fold)
    \item Overfitting warning if IS profitable but OOS negative
    \item Performance degradation percentage
\end{itemize}


% ##########################################################################
\chapter{Strategy Comparison and Anomaly Detection}
\label{ch:comparison}
% ##########################################################################

\section{Comparing Multiple Strategies}

The function \func{compare\_strategies} runs multiple strategies on the same data and returns a comparison DataFrame:

\begin{lstlisting}[caption={Strategy comparison usage.}]
result = compare_strategies(
    data,
    {"home": always_bet_home, "away": always_bet_away},
)
\end{lstlisting}

Returns a DataFrame with one row per strategy and columns for every metric in \func{calculate\_all\_metrics}.

\section{Ranking Strategies}

\func{rank\_strategies} sorts by a chosen metric and adds a \code{rank} column:

\begin{lstlisting}[caption={Ranking strategies.}]
ranked = rank_strategies(comparison_df, metric="sharpe_ratio")
\end{lstlisting}

\textbf{Dense ranking}: Tied strategies receive the same rank. If two strategies both have Sharpe = 1.5, they both get rank 1, and the next strategy gets rank 2.

\section{Display Comparison}

\func{display\_comparison} pretty-prints a formatted table with:
\begin{itemize}
    \item All numeric columns formatted to 4 decimal places
    \item ``Best by metric'' annotations (lowest for max\_drawdown, highest for everything else)
    \item Skips all-NaN columns to avoid crashes
    \item Skips ``Best by metric'' section if only one strategy is present
\end{itemize}

\section{Statistical Anomaly Detection}
\label{sec:anomaly-detection}

\func{detect\_suspicious\_results} provides six statistical checks that can catch cheating strategies which pass all 12 mechanical validator checks:

\begin{table}[H]
\centering
\caption{Anomaly detection checks.}
\label{tab:anomaly-checks}
\begin{tabularx}{\textwidth}{clX}
\toprule
\textbf{\#} & \textbf{Check} & \textbf{Description} \\
\midrule
1 & Win rate vs odds & Binomial test: is win rate significantly higher than implied probability? \\
2 & Perfect prediction & P(all wins by chance) under fair odds. Flagged if 100\% WR on $\geq 10$ trades. \\
3 & Odds-adjusted return & Actual PnL vs expected PnL from odds. Ratio $> 5\times$ is suspicious. \\
4 & Outcome entropy & Shannon entropy $< 0.2$ means outcomes are suspiciously predictable. \\
5 & Runs test & Tests randomness of W/L sequence. Too few runs = suspicious pattern. \\
6 & Bet-size correlation & Point-biserial correlation between bet size and outcome. $r > 0.3$ with $p < 0.05$ = suspicious. \\
\bottomrule
\end{tabularx}
\end{table}

\subsection{Check 1: Win Rate vs Implied Probability}

If the odds imply a 40\% win rate but the strategy wins 80\%, the binomial test quantifies how unlikely this is:

\begin{equation}
    p\text{-value} = P(X \geq \text{wins} \mid n, p_{\text{implied}})
\end{equation}

Flagged if $p < 0.001$.

\subsection{Check 3: Odds-Adjusted Return}

Expected PnL under fair pricing:
\begin{equation}
    \E[\pnl] = \sum_{i=1}^{N} \left[s_i(o_i - 1) \cdot \frac{1}{o_i} + (-s_i) \cdot \left(1 - \frac{1}{o_i}\right)\right] = 0
\end{equation}

At fair odds ($p_{\text{true}} = 1/o$), the expected PnL is exactly zero. Any large positive PnL at fair odds is suspicious.

When $\E[\pnl] \approx 0$, a ratio-based check would be unstable (division by near-zero). Instead, the check uses a per-trade excess threshold: profit $> 50\%$ of average bet size per trade at fair odds is flagged.

\subsection{Check 4: Shannon Entropy}
\begin{equation}
    H = -p_W \log_2(p_W) - p_L \log_2(p_L)
\end{equation}

Maximum entropy for binary outcomes is 1.0 (at $p_W = 0.5$). Entropy $< 0.2$ means outcomes are suspiciously predictable.

\subsection{Check 5: Wald-Wolfowitz Runs Test}

A ``run'' is a consecutive sequence of the same outcome (e.g., WWW or LL). Under randomness:
\begin{equation}
    \E[R] = 1 + \frac{2n_1 n_2}{n_1 + n_2}
\end{equation}
\begin{equation}
    \text{Var}(R) = \frac{2n_1 n_2 (2n_1 n_2 - n_1 - n_2)}{(n_1 + n_2)^2 (n_1 + n_2 - 1)}
\end{equation}

where $n_1$ = wins, $n_2$ = losses. Too few runs (all wins clustered together) suggests non-random outcomes.

\subsection{Check 6: Point-Biserial Correlation}

The \textbf{point-biserial correlation} between a binary variable (WIN/LOSS) and a continuous variable (bet size):

\begin{equation}
    r_{pb} = \frac{\bar{X}_1 - \bar{X}_0}{s_X} \sqrt{\frac{n_1 n_0}{n^2}}
\end{equation}

A strong positive correlation ($r > 0.3$, $p < 0.05$) means the strategy bets bigger on games it wins—which is exactly what a data-leaking strategy would do (it knows which games it will win and bets more on those).

\subsection{Overall Anomaly Score}

\begin{equation}
    \text{score} = \frac{\text{number of failed checks}}{\text{total checks run}}
\end{equation}

Flagged as suspicious if score $\geq 0.4$ or if win rate = 100\% on $\geq 10$ trades.


% ##########################################################################
\chapter{Testing}
\label{ch:testing}
% ##########################################################################

\section{Test Suite Overview}

The project has 200+ tests across 7 test files:

\begin{table}[H]
\centering
\caption{Test files and coverage areas.}
\label{tab:test-files}
\begin{tabularx}{\textwidth}{lrX}
\toprule
\textbf{File} & \textbf{Tests} & \textbf{Coverage} \\
\midrule
\file{test\_backtester\_backend.py}     & $\sim$25 & Core engine: basic, edge cases, cost model, Kelly \\
\file{test\_backtester\_extended.py}    & $\sim$25 & Extended: display, plotting, strategies, validation \\
\file{test\_audit\_fixes.py}           & 42+ & Audit: crash vectors, input validation, NaN, leakage, metrics \\
\file{test\_walk\_forward.py}          & $\sim$12 & Walk-forward: splits, folds, expanding window \\
\file{test\_comparison.py}             & 13  & Comparison, ranking, display, anomaly detection \\
\file{test\_strategies.py}             & $\sim$15 & Kelly criterion, strategy implementations \\
\file{test\_statistics.py}             & $\sim$10 & p-values, bootstrap, DSR, PBO \\
\bottomrule
\end{tabularx}
\end{table}

\section{Key Test Categories}

\subsection{Engine Crash Vectors}

Tests that the engine handles misbehaving strategies gracefully:
\begin{itemize}
    \item Strategy returns \code{None}, \code{"string"}, \code{[list]}, or \code{42} $\Rightarrow$ empty DataFrame returned
    \item Strategy raises exception $\Rightarrow$ prior trades preserved
    \item NaN bet size $\Rightarrow$ row skipped
\end{itemize}

\subsection{Input Validation Tests}

\begin{itemize}
    \item NaN initial\_bankroll $\Rightarrow$ \code{ValueError}
    \item cost\_pct $= 1.5$ $\Rightarrow$ \code{ValueError}
    \item cost\_flat $= -1$ $\Rightarrow$ \code{ValueError}
    \item Position sizing is case-insensitive: ``Kelly'' = ``kelly'' = ``KELLY''
    \item Kelly fraction out of range $\Rightarrow$ \code{ValueError}
\end{itemize}

\subsection{Data Leakage Tests}

\begin{itemize}
    \item \code{past\_games} does not contain \code{home\_win} column
    \item Modifying \code{context["history"]} does not affect next iteration
    \item Confidence values are clamped to $[0, 1]$
\end{itemize}

\subsection{Metric Known-Value Tests}

Tests with hand-calculated expected values:
\begin{itemize}
    \item Sharpe ratio with known returns (varying positive/negative)
    \item Sortino ratio with known downside
    \item Max drawdown with known equity curve
    \item Profit factor with known wins and losses
\end{itemize}


% ##########################################################################
\chapter{Complete Workflow Example}
\label{ch:workflow}
% ##########################################################################

This chapter walks through a complete end-to-end usage of the backtesting framework.

\section{Step 1: Load Data}

\begin{lstlisting}[caption={Loading historical game data.}]
from cuic_quant.backtest import load_backtest_data

data = load_backtest_data(
    start_date="2025-01-01",
    end_date="2025-12-31",
)
print(f"Loaded {len(data)} games")
\end{lstlisting}

\section{Step 2: Define a Strategy}

\begin{lstlisting}[caption={A simple value-betting strategy.}]
def value_bet_strategy(row, context=None):
    """Bet home when odds suggest value."""
    implied_prob = 1.0 / row["home_odds"]
    # Our model says home team is 10% more likely
    estimated_prob = implied_prob + 0.10

    if estimated_prob > implied_prob:
        return {
            "action": "BUY_HOME",
            "confidence": estimated_prob,
            "size": 100.0,
            "reason": f"Value: est={estimated_prob:.2f} > implied={implied_prob:.2f}",
        }
    return {"action": "SKIP", "confidence": 0, "size": 0}
\end{lstlisting}

\section{Step 3: Run Backtest}

\begin{lstlisting}[caption={Running the backtest.}]
from cuic_quant.backtest import backtest

results = backtest(
    data,
    value_bet_strategy,
    initial_bankroll=10000.0,
    cost_pct=0.02,        # 2% vig
    cost_flat=0.0,
    position_sizing="kelly",
    kelly_fraction=0.5,   # Half Kelly
)
print(f"Executed {len(results)} trades")
\end{lstlisting}

\section{Step 4: Validate Results}

\begin{lstlisting}[caption={Running the 12-check validator.}]
from cuic_quant.backtest import validate_backtest_results

validation = validate_backtest_results(results, data)
if validation["passed"]:
    print(f"All {validation['checks_run']} checks passed!")
else:
    for failure in validation["failures"]:
        print(f"FAIL: {failure}")
\end{lstlisting}

\section{Step 5: Compute Metrics}

\begin{lstlisting}[caption={Computing all performance metrics.}]
from cuic_quant.metrics import calculate_all_metrics

metrics = calculate_all_metrics(results)
print(f"Win rate:     {metrics['win_rate']:.1%}")
print(f"Total PnL:    ${metrics['total_pnl']:.2f}")
print(f"Sharpe ratio: {metrics['sharpe_ratio']:.4f}")
print(f"Max drawdown: {metrics['max_drawdown']:.1%}")
print(f"Kelly growth: {metrics['kelly_growth_rate']:.6f}")
\end{lstlisting}

\section{Step 6: Statistical Significance}

\begin{lstlisting}[caption={Testing statistical significance.}]
from cuic_quant.backtest.statistics import significance_report

report = significance_report(results)
print(f"P-value: {report['p_value']:.6f}")
print(f"Significant: {report['is_significant']}")
print(f"Sample size: {report['sample_size_assessment']}")
\end{lstlisting}

\section{Step 7: Walk-Forward Analysis}

\begin{lstlisting}[caption={Running walk-forward analysis.}]
from cuic_quant.backtest.walk_forward import (
    walk_forward_backtest,
    walk_forward_report,
)

wf = walk_forward_backtest(
    data, value_bet_strategy,
    n_splits=5,
    initial_bankroll=10000.0,
)
print(walk_forward_report(wf))
\end{lstlisting}

\section{Step 8: Compare Strategies}

\begin{lstlisting}[caption={Comparing multiple strategies.}]
from cuic_quant.backtest.comparison import (
    compare_strategies,
    rank_strategies,
    display_comparison,
)
from cuic_quant.backtest import always_bet_home, always_bet_away

comp = compare_strategies(data, {
    "value": value_bet_strategy,
    "home": always_bet_home,
    "away": always_bet_away,
})
ranked = rank_strategies(comp, metric="sharpe_ratio")
display_comparison(ranked)
\end{lstlisting}

\section{Step 9: Check for Data Leakage}

\begin{lstlisting}[caption={Statistical anomaly detection.}]
from cuic_quant.backtest.comparison import detect_suspicious_results

flags = detect_suspicious_results(results)
if flags["is_suspicious"]:
    print("WARNING: Results may indicate data leakage!")
    for check in flags["checks"]:
        status = "PASS" if check["passed"] else "FAIL"
        print(f"  [{status}] {check['name']}: {check['detail']}")
\end{lstlisting}


% ##########################################################################
\chapter{Decimal Odds Primer}
\label{ch:odds}
% ##########################################################################

Since the entire framework uses decimal odds, this chapter explains the conversion between odds formats.

\section{Decimal Odds}

Decimal odds represent the total return per unit wagered:
\begin{equation}
    \text{payout} = \text{stake} \times \text{decimal\_odds}
\end{equation}

\begin{example}
Odds of 2.50 mean: bet \$100, receive \$250 if you win (profit = \$150).
\end{example}

\section{Implied Probability}

\begin{equation}
    p_{\text{implied}} = \frac{1}{\text{odds}}
    \label{eq:implied-prob}
\end{equation}

\begin{table}[H]
\centering
\caption{Common decimal odds and their implied probabilities.}
\label{tab:odds-implied}
\begin{tabular}{ccc}
\toprule
\textbf{Decimal Odds} & \textbf{Implied Probability} & \textbf{Interpretation} \\
\midrule
1.50 & 66.7\% & Strong favourite \\
2.00 & 50.0\% & Even money \\
2.50 & 40.0\% & Slight underdog \\
3.00 & 33.3\% & Moderate underdog \\
5.00 & 20.0\% & Long shot \\
10.00 & 10.0\% & Very long shot \\
\bottomrule
\end{tabular}
\end{table}

\section{Bookmaker Margin (Vig/Overround)}

Bookmakers set odds so that implied probabilities sum to $> 100\%$:
\begin{equation}
    \text{overround} = \frac{1}{o_H} + \frac{1}{o_A} - 1
\end{equation}

\begin{example}
Home odds 1.90, away odds 2.00:
\[
    \text{overround} = \frac{1}{1.90} + \frac{1}{2.00} = 0.5263 + 0.5000 = 1.0263
\]
The bookmaker has a 2.63\% margin.
\end{example}

\section{Converting From American Odds}

\begin{align}
    \text{American} > 0: \quad \text{decimal} &= \frac{\text{American}}{100} + 1 \\
    \text{American} < 0: \quad \text{decimal} &= \frac{100}{|\text{American}|} + 1
\end{align}


% ##########################################################################
\chapter{Error Handling and Edge Cases}
\label{ch:edge-cases}
% ##########################################################################

\section{Engine Edge Cases}

\begin{table}[H]
\centering
\caption{How the engine handles edge cases.}
\label{tab:edge-cases}
\begin{tabularx}{\textwidth}{lX}
\toprule
\textbf{Edge Case} & \textbf{Behaviour} \\
\midrule
Empty input data & Returns empty DataFrame with correct 11 columns \\
All rows skipped & Returns empty DataFrame with correct 11 columns \\
Bankroll reaches zero & Loop breaks, returns trades executed so far \\
NaN odds & Row skipped silently \\
Odds $\leq 1.0$ & Row skipped silently (invalid decimal odds) \\
NaN home\_win & Row skipped with warning \\
home\_win $\notin \{0, 1\}$ & Row skipped with warning \\
String-typed odds & Converted via \code{float()}; if that fails, row skipped with warning \\
Strategy returns \code{None} & Warning emitted, row skipped \\
Strategy raises exception & Warning emitted, row skipped, prior trades preserved \\
NaN bet size from Kelly & Row skipped with warning \\
cost\_flat $>$ bankroll & Bet size becomes 0, row skipped \\
\bottomrule
\end{tabularx}
\end{table}

\section{Metrics Edge Cases}

\begin{table}[H]
\centering
\caption{How metrics handle edge cases.}
\label{tab:metrics-edge}
\begin{tabularx}{\textwidth}{lX}
\toprule
\textbf{Edge Case} & \textbf{Behaviour} \\
\midrule
Zero trades & All metrics return 0.0 \\
One trade & Sharpe returns 0.0 (need $\geq 2$ for std dev) \\
All wins, no losses & Profit factor returns $\infty$ \\
All losses, no wins & Profit factor returns 0.0 \\
Constant returns (zero variance) & Sharpe returns 0.0 \\
No downside (all positive returns) & Sortino returns 0.0 \\
No confidence values & Brier Score and Log Loss return NaN \\
No closing odds & CLV returns NaN \\
Timestamps unavailable & Periods per year defaults to 365 with warning \\
\bottomrule
\end{tabularx}
\end{table}


% ##########################################################################
\chapter{Module Reference}
\label{ch:reference}
% ##########################################################################

This chapter provides a quick-reference listing of every public function in the framework.

\section{backtest.engine}

\begin{longtable}{lp{9cm}}
\toprule
\textbf{Function} & \textbf{Description} \\
\midrule
\endhead
\func{backtest} & Core backtest loop. Takes data + strategy, returns 11-column DataFrame. \\
\func{always\_bet\_home} & Reference strategy: always bets \$100 on home team. \\
\func{always\_bet\_away} & Reference strategy: always bets \$100 on away team. \\
\func{kelly\_bet\_home\_demo} & Demo strategy showing Kelly plumbing (not a real edge model). \\
\bottomrule
\end{longtable}

\section{backtest.data\_loader}

\begin{longtable}{lp{9cm}}
\toprule
\textbf{Function} & \textbf{Description} \\
\midrule
\endhead
\func{load\_backtest\_data} & Load games from Railway DB or CSV. Returns sorted DataFrame. \\
\bottomrule
\end{longtable}

\section{backtest.validator}

\begin{longtable}{lp{9cm}}
\toprule
\textbf{Function} & \textbf{Description} \\
\midrule
\endhead
\func{validate\_backtest\_results} & 12-check validation (schema, math, leakage). Auto-reads cost params from \code{.attrs}. \\
\bottomrule
\end{longtable}

\section{metrics}

\begin{longtable}{lp{9cm}}
\toprule
\textbf{Function} & \textbf{Description} \\
\midrule
\endhead
\func{calculate\_all\_metrics} & Compute all metrics from trade DataFrame. \\
\func{calculate\_sharpe\_ratio} & Annualised Sharpe ratio. \\
\func{calculate\_sortino\_ratio} & Annualised Sortino ratio (downside deviation). \\
\func{calculate\_max\_drawdown} & Maximum drawdown as fraction of peak equity. \\
\func{calculate\_win\_rate} & Win rate from WIN/LOSS outcomes. \\
\func{calculate\_profit\_factor} & Gross profit / gross loss. \\
\func{calculate\_kelly\_growth\_rate} & Realised log growth rate per bet. \\
\func{calculate\_clv} & Average Closing Line Value. \\
\func{calculate\_brier\_score} & Brier Score for probability calibration. \\
\func{calculate\_log\_loss} & Log loss (cross-entropy). \\
\bottomrule
\end{longtable}

\section{backtest.statistics}

\begin{longtable}{lp{9cm}}
\toprule
\textbf{Function} & \textbf{Description} \\
\midrule
\endhead
\func{calculate\_p\_value} & Binomial test p-value (one-sided by default). \\
\func{bootstrap\_confidence\_interval} & Non-parametric bootstrap CI for any metric. \\
\func{minimum\_sample\_size} & Min bets needed for significance. \\
\func{significance\_report} & Full significance report with p-value, CI, sample size. \\
\func{bonferroni\_correction} & Bonferroni multiple-testing correction. \\
\func{holm\_bonferroni\_correction} & Holm-Bonferroni step-down procedure. \\
\func{deflated\_sharpe\_ratio} & DSR (Bailey \& LdP 2014) with skewness/kurtosis. \\
\func{probability\_of\_backtest\_overfitting} & Simplified PBO from single IS/OOS split. \\
\func{overfitting\_report} & Full overfitting report for multiple strategies. \\
\bottomrule
\end{longtable}

\section{backtest.walk\_forward}

\begin{longtable}{lp{9cm}}
\toprule
\textbf{Function} & \textbf{Description} \\
\midrule
\endhead
\func{train\_test\_split} & Simple chronological split with optional purge gap. \\
\func{walk\_forward\_backtest} & Rolling-window walk-forward with $n$ folds. \\
\func{expanding\_window\_backtest} & Growing training window walk-forward. \\
\func{anchored\_walk\_forward} & Walk-forward anchored to a fixed start date. \\
\func{combinatorial\_purged\_cv} & CPCV from Lopez de Prado. \\
\func{walk\_forward\_report} & Human-readable walk-forward report. \\
\bottomrule
\end{longtable}

\section{backtest.comparison}

\begin{longtable}{lp{9cm}}
\toprule
\textbf{Function} & \textbf{Description} \\
\midrule
\endhead
\func{compare\_strategies} & Run multiple strategies, return comparison DataFrame. \\
\func{rank\_strategies} & Sort and rank by chosen metric (dense ranking). \\
\func{display\_comparison} & Pretty-print comparison table. \\
\func{detect\_suspicious\_results} & 6-check statistical anomaly detection. \\
\func{implied\_edge\_analysis} & Actual vs implied win rate and edge. \\
\func{validate\_no\_leakage\_statistical} & Statistical complement to mechanical validator. \\
\bottomrule
\end{longtable}

\section{strategies.kelly\_criterion}

\begin{longtable}{lp{9cm}}
\toprule
\textbf{Function} & \textbf{Description} \\
\midrule
\endhead
\func{calculate\_kelly\_fraction} & Optimal Kelly bet fraction with max cap. \\
\func{calculate\_kelly\_for\_multiple\_bets} & Kelly for simultaneous uncorrelated bets. \\
\func{expected\_growth\_rate} & Expected log growth rate for a given bet fraction. \\
\func{kelly\_from\_american\_odds} & Kelly with American odds input. \\
\bottomrule
\end{longtable}


% ##########################################################################
\chapter{Glossary}
\label{ch:glossary}
% ##########################################################################

\begin{description}[style=nextline, leftmargin=3cm]
    \item[Bankroll] The total capital available for betting. Starts at \code{initial\_bankroll}.
    \item[BUY\_HOME / BUY\_AWAY] Actions indicating which side to bet on.
    \item[CLV] Closing Line Value — the gold-standard metric for betting skill.
    \item[CPCV] Combinatorial Purged Cross-Validation — tests all possible train/test partitions.
    \item[Data Leakage] When a strategy has access to information it wouldn't have in live betting (e.g., the game outcome).
    \item[Decimal Odds] Total payout per unit wagered (e.g., 2.50 means \$2.50 returned per \$1.00 bet).
    \item[DSR] Deflated Sharpe Ratio — adjusts for multiple testing.
    \item[Edge] The difference between estimated and implied win probability.
    \item[Fold] A single train/test partition in walk-forward analysis.
    \item[Half Kelly] Betting half the Kelly-optimal fraction for reduced variance.
    \item[Kelly Criterion] Formula for optimal bet sizing that maximises log growth rate.
    \item[OOS] Out-of-sample — data not used for training/fitting.
    \item[Overround / Vig] The bookmaker's margin (implied probabilities sum to $> 100\%$).
    \item[PBO] Probability of Backtest Overfitting.
    \item[PnL] Profit and Loss — the dollar gain or loss from a trade.
    \item[Purge Gap] Rows removed between train and test sets to prevent leakage from look-back features.
    \item[Sharpe Ratio] Risk-adjusted return: mean excess return divided by standard deviation.
    \item[Signal] The dictionary returned by a strategy function (action, confidence, size, reason).
    \item[SKIP] Strategy action indicating no bet should be placed.
    \item[Sortino Ratio] Like Sharpe but only penalises downside volatility.
    \item[Walk-Forward] Testing methodology where the model is trained on past data and tested on future data.
\end{description}


% ##########################################################################
% APPENDIX
% ##########################################################################
\appendix

\chapter{Mathematical Derivations}
\label{app:derivations}

\section{Sharpe Ratio Annualisation}

Starting from per-period returns $r_1, r_2, \ldots, r_N$ with mean $\mu$ and standard deviation $\sigma$:

If there are $P$ periods per year, and assuming returns are i.i.d.:
\begin{align}
    \mu_{\text{annual}} &= P \cdot \mu \\
    \sigma_{\text{annual}} &= \sqrt{P} \cdot \sigma
\end{align}

Therefore:
\begin{equation}
    \text{SR}_{\text{annual}} = \frac{\mu_{\text{annual}}}{\sigma_{\text{annual}}} = \frac{P\mu}{\sqrt{P}\sigma} = \frac{\mu}{\sigma}\sqrt{P}
\end{equation}

\section{Downside Deviation Derivation}

The downside deviation (Sortino \& Price 1994) is defined as:
\begin{equation}
    \sigma_d = \sqrt{\E[\min(r - r_f, 0)^2]}
\end{equation}

The sample estimator uses all $N$ observations:
\begin{equation}
    \hat{\sigma}_d = \sqrt{\frac{1}{N}\sum_{i=1}^{N}[\min(r_i - r_f, 0)]^2}
\end{equation}

Note: some implementations use only the negative returns in the denominator ($N_d$ instead of $N$). Our implementation uses $N$ (all observations) which is the standard academic definition and avoids overstating downside risk.

\section{Kelly Criterion Optimality}

The Kelly Criterion maximises the expected geometric growth rate:
\begin{equation}
    G(f) = \E[\ln(W_{t+1}/W_t)] = p\ln(1+bf) + q\ln(1-f)
\end{equation}

First-order condition:
\begin{equation}
    G'(f^*) = \frac{pb}{1+bf^*} - \frac{q}{1-f^*} = 0
\end{equation}

Solving:
\begin{equation}
    f^* = \frac{bp - q}{b} = p - \frac{q}{b} = p - \frac{1-p}{o-1}
\end{equation}

Second-order condition (concavity):
\begin{equation}
    G''(f) = -\frac{pb^2}{(1+bf)^2} - \frac{q}{(1-f)^2} < 0 \quad \forall f \in (0,1)
\end{equation}

This confirms $f^*$ is indeed a maximum.

\section{Expected Maximum Sharpe Ratio}

For $k$ independent strategies with Sharpe ratios drawn from $\mathcal{N}(0, \sigma_{\text{SR}}^2)$, the expected maximum follows from extreme value theory:

\begin{equation}
    \E[\max(SR_1, \ldots, SR_k)] \approx \sigma_{\text{SR}}\left[(1-\gamma)\Phi^{-1}\!\left(1-\frac{1}{k}\right) + \gamma\Phi^{-1}\!\left(1-\frac{1}{ke}\right)\right]
\end{equation}

where $\gamma \approx 0.5772$ is the Euler-Mascheroni constant.

This grows approximately as $\sigma_{\text{SR}} \sqrt{2\ln k}$ for large $k$.

\chapter{File Structure}
\label{app:file-structure}

\begin{lstlisting}[language={}, basicstyle=\ttfamily\small]
CUIC_Sem2_Project/
|-- src/cuic_quant/
|   |-- backtest/
|   |   |-- __init__.py              # Re-exports all public functions
|   |   |-- backtester_backend.py    # Facade (re-exports from submodules)
|   |   |-- engine.py                # Core backtest loop
|   |   |-- data_loader.py           # DB/CSV data loading
|   |   |-- validator.py             # 12-check validation
|   |   |-- walk_forward.py          # Walk-forward analysis (6 functions)
|   |   |-- comparison.py            # Strategy comparison + anomaly detection
|   |   |-- statistics.py            # Statistical significance + overfitting
|   |   \-- visualization.py         # Plotting functions
|   |-- metrics/
|   |   \-- __init__.py              # All metric calculations
|   \-- strategies/
|       \-- kelly_criterion.py       # Kelly Criterion implementation
|-- tests/
|   |-- test_backtester_backend.py   # Core engine tests
|   |-- test_backtester_extended.py  # Extended feature tests
|   |-- test_audit_fixes.py          # Comprehensive audit fix tests
|   |-- test_walk_forward.py         # Walk-forward tests
|   |-- test_comparison.py           # Comparison + anomaly tests
|   |-- test_strategies.py           # Kelly + strategy tests
|   \-- test_statistics.py           # Statistical tests
|-- data/
|   \-- dummy_backtest_input.csv     # 25-row synthetic NBA data
|-- tools/
|   \-- backtester.ipynb             # Interactive notebook
\-- docs/
    \-- backtester_reference.tex     # This document
\end{lstlisting}


% ##########################################################################
% REFERENCES
% ##########################################################################

\chapter*{References}
\addcontentsline{toc}{chapter}{References}

\begin{enumerate}[label={[\arabic*]}]
    \item Bailey, D.H. \& Lopez de Prado, M. (2014). ``The Deflated Sharpe Ratio: Correcting for Selection Bias, Backtest Overfitting and Non-Normality.'' \textit{Journal of Portfolio Management}, 40(5), 94--107.

    \item Bailey, D.H., Borwein, J.M., Lopez de Prado, M. \& Zhu, Q.J. (2017). ``The Probability of Backtest Overfitting.'' \textit{Journal of Computational Finance}, 20(4), 39--69.

    \item Kelly, J.L. (1956). ``A New Interpretation of Information Rate.'' \textit{Bell System Technical Journal}, 35(4), 917--926.

    \item Thorp, E.O. (2006). ``The Kelly Criterion in Blackjack, Sports Betting, and the Stock Market.'' In \textit{Handbook of Asset and Liability Management}, Volume 1.

    \item Sortino, F.A. \& Price, L.N. (1994). ``Performance Measurement in a Downside Risk Framework.'' \textit{Journal of Investing}, 3(3), 59--64.

    \item Sharpe, W.F. (1966). ``Mutual Fund Performance.'' \textit{Journal of Business}, 39(1), 119--138.

    \item Lopez de Prado, M. (2018). \textit{Advances in Financial Machine Learning}. Wiley.
\end{enumerate}


% ==========================================================================
\end{document}
% ==========================================================================
